\documentclass[12pt]{article}

% The preamble file contains all the macros and page setup. I recommend making a folder where you compile your exercises and always call on the preamble.
%%%%%%%%%%%%%%%%%%%%%%%%%%%%%%%%%%%%%%%%%%%%%%%%%%%%%%%%%%%%%%%%%%%%%%
%% all the formatting stuff and packages



\usepackage{geometry}
\geometry{
  top=4cm,
  headheight=70pt,
  headsep=25pt,
  inner=1.25in,
  outer=1.25in
}

\usepackage[utf8]{inputenc}

\usepackage{xcolor,graphicx,hyperref,amsthm,amsmath,amssymb}
\usepackage{algpseudocode}


\usepackage{fancyhdr}
\pagestyle{fancy}

%%% Begin Header Definitions %%%
\fancypagestyle{firststyle}
{
   	\lhead{}
	\rhead{\raisebox{0.2cm}{\includegraphics[width=3.5cm]{resource/hpi_logo.pdf}}}
	\renewcommand{\headrulewidth}{0pt}
}


\newcommand{\makeheader}[1]{
\noindent
\textbf{Algorithmix \hfill Winter term 2019/20 \\ Andreas G\"obel}

\begin{center}
\large{\textbf{\makeHeadline}}\\
\end{center}

\begin{center}
\textit{#1} 
\end{center}


}

\newcommand{\lecturer}[1]{\newcommand{\makeLecturer}{#1}}
\newcommand{\lectureName}[1]{\newcommand{\makeLectureName}{#1}}
\newcommand{\semester}[1]{\newcommand{\makeSemester}{#1}}
\newcommand{\headline}[1]{\newcommand{\makeHeadline}{#1}}
\newcommand{\urlLine}[1]{\newcommand{\makeUrlLine}{#1}}
%%% End Header Definitions %%%

\newcommand{\pr}{\ensuremath{\mathbb{P}}}
\newcommand{\E}{\ensuremath{\mathbb{E}}}


\newcommand{\makeCompact}{%
\setlength{\itemsep}{1pt}%
\setlength{\parsep}{0pt}%
\setlength{\parskip}{0pt}}


\newcommand{\Var}{\mathrm{Var}}


\newcommand{\student}[2]{
\thispagestyle{firststyle}

%%% Begin Header Specification %%%

\headline{Solutions to homework #1}

\makeheader{#2}

\medskip

}

\newcommand{\sol}[1]{ \hspace{.5cm}\\ \noindent \textbf{Exercise #1.}}

% Here you can define your own macros
\newcommand{\calF}{\ensuremath{\mathcal{F}}}

\begin{document}
 
 %To create a header for the document use the following macro. \student{Homework number}{Name}

 \student{2}{Stefan Reschke}

 \sol{1} (a) We proof given cases of the definition. \begin{enumerate}
 	\item If $a$ and $b$ are even numbers, we have \begin{align}
 		\texttt{gcd}(a,b) &= \texttt{gcd}\bigg(2\cdot\frac{a}{2},2\cdot\frac{b}{2}\bigg) \\
 		&=2\cdot\texttt{gcd}\bigg(\frac{a}{2},\dfrac{b}{2}\bigg),
 	\end{align} which works, because $\frac{a}{2},\frac{b}{2}\in\mathbb{N}$.
 	\item If $a$ is odd and $b$ is even, we have \begin{align}
 		\texttt{gcd}(a,b)&=\texttt{gcd}\bigg(a,2\cdot\frac{b}{2}\bigg) \\
 		&=\texttt{gcd}\bigg(a,\frac{b}{2}\bigg)\cdot\texttt{gcd}(a,2).
 	\end{align} Since $a$ is odd, we know that $\texttt{gcd}(a,2)=1$. Thus \begin{align*}
 		\texttt{gcd}(a,b)&=\texttt{gcd}\bigg(a,\frac{b}{2}\bigg).
 	\end{align*}
 \end{enumerate}
\end{document}
